\centerline{\Huge Условия и решения задач.}
\centerline{\bf \Huge 5 класс}

\title{Условия и решения задач}

\problem{5.1} Найдите наибольшее и наименьшее девятизначное число, которое делится на $2023$. \par
\textbf{Решение.} Разделим с остатком $100 \ 000 \ 000$ на $2023$ и получим, что $100 \ 000 \ 000 = 2023\cdot 49431 + 1087$. Тогда искомое наименьшее девятизначное число это $100 \ 000 \ 000 + (2023 - 1087) = 100 \ 000 936.$ \par
Точно так же найдем наибольшее. Разделим с остатком $999 \ 999 \ 999$ на $2023$ и получим, что $999 \ 999 \ 999 = 2023\cdot 494315 + 754$. Тогда искомое наименьшее девятизначное число это $999 \ 999 \ 999 - 754 = 999 \ 999 \ 245.$ \\

\textbf{Критерии.} \\
\textit{7 баллов} $-$ Полностью верное решение. \\
\textit{5 баллов} $-$ Верно выражены оба искомых числа, но одно вычислено с ошибкой, а другое нет. \\ 
\textit{3 балла} $-$ Верно выражены оба искомых числа, но оба вычислены с ошибкой. \\ 
\textit{0-2 балла} $-$ Есть какие-то продвижения, которые могут привести к решению. \\ 
\textit{0 баллов} $-$ Решение отсутствует. Решение неверное, продвижения отсутствуют. \\ 

\problem{5.2} В Калмыкии между любыми двумя населенными пунктами есть ровно одна дорога. Докажите, что найдутся два населенных пункта, из которых исходит одинаковое количество дорог. \par

\textbf{Решение.} Пусть в Калмыкии $n$ населённых пунктов. Из каждого населенного пункта может исходить либо 0, либо 1, либо 2, $\dotsc$, либо $n-1$ дорога. Заметим, что не может одновременно существовать два населённых пункта из которого исходит 0 и из которого исходит $n-1$ дорога. Так как первый не будет соединен ни с кем, а второй должен быть соединен со всеми, а значит и с первым тоже. То есть из этих двух вариантов может существовать только один. Тогда у каждого населенного пункта $n-1$ вариант, сколько из него будет исходить дорог. Населенных пунктов у нас $n$, а количество вариантов исходящих из них дорог $n-1$. Значит по принципу Дирихле есть два пункта, их которых исходит одинаковое число дорог. 
\\

\textbf{Критерии.} \\
\textit{7 баллов} $-$ Полностью верное решение. \newpage
\textit{4 балла} $-$ Доказано, что не может существовать пунктов с исходящими 0 и $n-1$ дорогами одновременно.\\ 
\textit{2 балла} $-$ Написано, что из каждого населенного пункта может исходидить либо 0, либо 1, $\dotsc$, либо $n-1$ дорога. \\
\textit{0-2 балла} $-$ Есть какие-то продвижения, которые могут привести к решению. \\ 
\textit{0 баллов} $-$ Решение отсутствует. Решение неверное, продвижения отсутствуют. \\ 

\problem{5.3} На острове Парадиз проживают $N$ островитян: эльдийцы, которые всегда говорят правду и марлийцы, которые всегда лгут. Путешественник задал им вопрос: «Кого на острове больше $-$ эльдийцев или марлийцев?» \\
Первый житель острова ответил «Марлийцев больше» \\ 
Второй $-$ «Эльдийцев не меньше» \\
Третий «Марлийцев больше» \\
Четвертый «Эльдийцев не меньше» \\
… \\
$(N-1)-$ый ответил «Эльдийцев не меньше» \\
$N-$ый «Марлийцев больше». \\
Может ли $N$ равняться $2023$? \par

\textbf{Решение.} Заметим, что высказывание «Марлийцев больше» противоположно высказыванию «Эльдийцев не меньше». Стало быть одно из них правда, а другое ложь. Следовательно среди отвечавших марлийцы и эльдийцы чередовались. Предположим, что $N = 2023$. Тогда высказываний «Марлийцев больше» было произнесено $1012$ раз, а «Эльдийцев не меньше» $1011$ раз. Пусть «Марлийцев больше» говорили эльдийцы, которые всегда говорят правду. Но тогда эльдийцев $1012$ и они соврали, но они всегда говорят правду. Противоречие. Пусть «Марлийцев больше» говорили марлийцы, которые всегда лгут. Тогда марлийцев $1012$ и они сказали правду, но они всегда лгут. Противоречие. Следовательно $N$ не может равняться $2023$. \\

\textbf{Критерии.} \\
\textit{7 баллов} $-$ Полностью верное решение.\\ 
\textit{4 балла} $-$ Понимание того, что высказывания чередуются среди марлийцев и эльдийцев.\\
\textit{3 балла} $-$ Доказано, что высказывания «Марлийцев больше» и «Эльдийцев не меньше» противоречивы.\\
\textit{0-2 балла} $-$ Есть какие-то продвижения, которые могут привести к решению. \\
\textit{0 баллов} $-$ Решение отсутствует. Решение неверное, продвижения.\\

\problem{5.4} У Аркадия есть клетчатая доска $20\times23$, в каждой клетке которой сидит по кузнечику. В какой-то момент времени каждый кузнечик перепрыгнул на соседнюю по вершине клетку. Докажите, что после этого найдутся хотя бы 20 клеток доски, которые будут пустыми. \par
\textbf{Решение.} Не нарушая общности, пусть у Аркадия была доска, состоящая из 20 строк и 23 столбцов. Покрасим первый столбец в красный, второй в синий, третий в красный, четвертый в синий, $\dotsc$, $22-$ой в синий, $23-$й в красный. Тогда красных клеток у нас $12\cdot20 = 240$, а кузнечиков сидящих на синих клетках $11\cdot20 = 220$. Заметим, что кузнечик сидящий на красной клетке, после прыжка будет сидеть на синей и наоборот. Следовательно 220 кузнечиков перепрыгнут на 240 красных клеток. Но даже если все 220 кузнечиков сядут на разные клетки, то свободных красных будет хотя бы 20. Что и требовалось доказать. \\

\textbf{Критерии.} \\
\textit{7 баллов} $-$ Полностью верное решение.\\ 
\textit{4 балла} $-$ Доказано, что кузнечник меняет цвет клетки при прыжке.\\ 
\textit{2 балла} $-$ Правильно выбранная раскраска.\\
\textit{0-2 балла} $-$ Есть какие-то продвижения, которые могут привести к решению.\\
\textit{0 баллов} $-$ Решение отсутствует. Решение неверное, продвижения отсутствуют. \\ 


\problem{5.5} На бесконечной доске выписаны числа $5, \ 55, \ 555, \ \dots  \ , \ \underbrace{555\dots 5}_{n}$, \ \dots \\
Докажите, что среди выписанных найдутся два, разность которых делится на $2023$. 
\par
\textbf{Решение.} Рассмотрим остатки этих чисел при делении на $2023$. Заметим, что любое число может давать остатки 0, 1, 2, \dots , $2022$ при делении на $2023$. Так как у нас бесконечное количество чисел(больше $2023$ в том числе), то по принципу Дирихле найдутся два числа имеющие один и тот же остаток при делении на $2023$. Разность этих чисел и будет делиться на $2023$.

\textbf{Критерии.} \\
\textit{7 баллов} $-$ Полностью верное решение. \\
\textit{5 баллов} $-$ Доказано, что найдутся два числа с одним и тем же остатком при делении на 2023.\\
\textit{0-2 балла} $-$ Есть какие-то продвижения, которые могут привести к решению. \\ 
\textit{0 баллов} $-$ Решение отсутствует. Решение неверное, продвижения отсутствуют. \\ 